% ===== PRÉAMBULE CANONIQUE — à coller VERBATIM en tête de CHAQUE .tex =====
\documentclass[11pt,a4paper]{article}
\usepackage[utf8]{inputenc}\usepackage[T1]{fontenc}\usepackage{lmodern}
\usepackage[french]{babel}
\usepackage{amsmath,amssymb,mathtools}\usepackage{icomma}
\usepackage[version=4]{mhchem}          % \ce{...} : équations & formules
\usepackage{siunitx}\sisetup{locale=FR} % \qty{}{}, \si{}, \ang{}
\usepackage{chemfig}                    % \chemfig{...} : structures développées
\usepackage{xcolor}\usepackage{colortbl}
\usepackage{tikz}\usepackage{pgfplots}\pgfplotsset{compat=1.18}
\usepackage[most]{tcolorbox}\usepackage{enumitem}\usepackage{array,booktabs}
\usepackage[a4paper,margin=2.2cm]{geometry}
\usetikzlibrary{arrows.meta,calc,angles,quotes,positioning,patterns,backgrounds,shapes.geometric}
\definecolor{primary}{RGB}{222,49,99}\definecolor{primarydark}{RGB}{150,22,55}
\definecolor{defcol}{RGB}{0,114,178}\definecolor{methcol}{RGB}{52,92,148}
\definecolor{excol}{RGB}{0,135,90}\definecolor{attncol}{RGB}{200,40,55}
\tcbset{boxbase/.style={enhanced,breakable,boxrule=0.9pt,arc=2.5pt,
  left=3mm,right=3mm,top=2.3mm,bottom=2.3mm,fonttitle=\bfseries,coltitle=white}}
% --- boîtes TUTORIEL (noms EXACTS, ne pas renommer) ---
\newtcolorbox{cle}[1][]{boxbase,colback=defcol!6,colframe=defcol,
  title={\textsc{À retenir}\ifx\relax#1\relax\relax\else\textmd{ : \itshape #1}\fi}}
\newtcolorbox{methode}[1][]{boxbase,colback=methcol!7,colframe=methcol,sharp corners,
  title={\textsc{Méthode}\ifx\relax#1\relax\relax\else\textmd{ --- \itshape #1}\fi}}
\newtcolorbox{exemplebox}[1][]{boxbase,colback=excol!5,colframe=excol,
  title={\textsc{Exemple}\ifx\relax#1\relax\relax\else\textmd{ : \itshape #1}\fi}}
\newtcolorbox{piege}[1][]{boxbase,colback=attncol!6,colframe=attncol,
  title={\textsc{Piège}\ifx\relax#1\relax\relax\else\textmd{ : \itshape #1}\fi}}
\newtcolorbox{remarque}[1][]{boxbase,colback=black!4,colframe=black!45,
  title={\textsc{Remarque}\ifx\relax#1\relax\relax\else\textmd{ : \itshape #1}\fi}}
% --- boîtes EXERCICE / CORRIGÉ (noms EXACTS, auto-comptées) ---
\newtcolorbox[auto counter]{exo}[1][]{boxbase,colback=primary!4,colframe=primary,
  title={\textsc{Exercice}~\thetcbcounter\ifx\relax#1\relax\relax\else\textmd{ --- \itshape #1}\fi}}
\newtcolorbox[auto counter]{corrige}[1][]{boxbase,colback=excol!4,colframe=excol,
  title={\textsc{Corrigé}~\thetcbcounter\ifx\relax#1\relax\relax\else\textmd{ --- \itshape #1}\fi}}
\newcommand{\R}{\mathbb{R}}\newcommand{\N}{\mathbb{N}}\newcommand{\Z}{\mathbb{Z}}\newcommand{\ds}{\displaystyle}
% (un \usepackage{fancyhdr}+en-tête est possible ; il est ignoré côté web)
% =========================================================================
\begin{document}
% THEME_TITLE: Configuration électronique en couches et électrons de valence

Les électrons se répartissent autour du noyau sur des \textbf{couches} successives. Savoir les remplir
donne la \textbf{configuration électronique}, d'où l'on lit les \textbf{électrons de valence} qui
gouvernent toute la chimie de l'atome.

\begin{cle}[les couches et leur capacité]
Les couches se notent K, L, M, N\dots\ (numérotées $n = 1, 2, 3, 4\dots$). La couche $n$ contient au
maximum $\mathbf{2n^2}$ électrons :
\[ \text{K}:2,\qquad \text{L}:8,\qquad \text{M}:18,\qquad \text{N}:32. \]
On remplit \textbf{de l'intérieur vers l'extérieur} (K, puis L, puis M\dots). Dans ce modèle simplifié,
la \textbf{couche la plus externe ne dépasse jamais $8$ électrons} : dès qu'elle atteint $8$, on ouvre
la couche suivante.
\end{cle}

\begin{cle}[configuration et électrons de valence]
La \textbf{configuration électronique} s'écrit $(\mathrm{K})^a(\mathrm{L})^b(\mathrm{M})^c\dots$, où les
exposants comptent les électrons de chaque couche ; leur somme vaut $n(e^-)$ (donc $Z$ pour un atome
neutre). La \textbf{couche de valence} est la \emph{dernière} couche occupée, et ses électrons sont les
\textbf{électrons de valence}.
\end{cle}

\begin{methode}[établir la configuration d'un atome neutre]
\begin{enumerate}[leftmargin=1.7em,topsep=2pt,itemsep=1pt]
  \item compter les électrons : $n(e^-) = Z$ ;
  \item remplir K ($2$ au plus), puis L ($8$), puis M\dots, en gardant la couche externe $\le 8$ ;
  \item la dernière couche occupée est la couche de valence ; ses électrons sont les électrons de valence.
\end{enumerate}
Pour un \textbf{ion} : on part de l'atome neutre, puis on \emph{retire} (cation) ou \emph{ajoute}
(anion) le bon nombre d'électrons.
\end{methode}

\begin{exemplebox}[le phosphore $Z = 15$]
$15$ électrons : K$^2$, puis L$^8$ ($2+8=10$), puis les $5$ restants sur M. D'où
$(\mathrm{K})^2(\mathrm{L})^8(\mathrm{M})^5$ : la couche de valence est $\mathbf{(M)^5}$, avec $5$
électrons de valence.
\end{exemplebox}

\begin{piege}[la couche de valence est la dernière occupée]
Pour le phosphore, la valence est $(\mathrm{M})^5$, \textbf{pas} $(\mathrm{L})^5$ : la couche L est
\emph{saturée} à $8$. On ne « saute » jamais une couche interne pour désigner la valence.
\end{piege}

\begin{remarque}[la couche M : tantôt $8$, tantôt $18$]
Tant que M est la couche \emph{externe}, elle est plafonnée à $8$ (argon : $(\mathrm{K})^2(\mathrm{L})^8
(\mathrm{M})^8$). Dès qu'une couche N s'ouvre au-dessus, M devient \emph{interne} et peut se remplir
jusqu'à $18$ (brome : $(\mathrm{K})^2(\mathrm{L})^8(\mathrm{M})^{18}(\mathrm{N})^7$).
\end{remarque}

\begin{cle}[symbole de Lewis]
On dispose les \textbf{électrons de valence} en points autour du symbole, un par côté d'abord, puis on
les \textbf{apparie} : les paires sont des \textbf{doublets}, les points restés seuls des
\textbf{électrons célibataires}. Exemple : l'oxygène ($6$ électrons de valence) $\to$ $2$ doublets et
$2$ célibataires.
\end{cle}

\end{document}
